\subsection{Node, Link ও Port-এর হিসাব: শর্তসহ ছক}
\begin{tabularx}{\linewidth}{|p{.17\linewidth}|p{.30\linewidth}|X|}\hline
\textbf{বিন্যাস} & \textbf{Link-এর সংখ্যা} & \textbf{শর্ত ও Port-এর অর্থ}\\\hline
Star & $n$ প্রান্তিক ডিভাইস হলে $n$; কেন্দ্রসহ মোট $N$ হলে $N-1$ & প্রতিটি প্রান্তিক ডিভাইসের জন্য একটি কেন্দ্রীয় লিংক; সুইচে অন্তত $n$ পোর্ট। ৭ PC + ১ সুইচ = ৭ লিংক, ৮ ডিভাইস।\\\hline
Ring & $n$ নোডে $n$ লিংক & একটি সাধারণ বন্ধ চক্রে প্রত্যেক নোডের দুই প্রতিবেশী থাকে। এটি মোট সংযোগের সংখ্যা; নির্দিষ্ট বার্তার পথের ধাপসংখ্যা নয়।\\\hline
Tree & মোট $N$ শীর্ষবিন্দুতে $N-1$ সংযোগরেখা & সংযুক্ত, চক্রবিহীন গাছের বিন্যাস। কেন্দ্রীয় ও মধ্যবর্তী ডিভাইসও নোড হিসেবে গণ্য। বাড়তি বিকল্প লিংক থাকলে সরল সূত্রটি প্রযোজ্য নয়।\\\hline
Full Mesh & $L=n(n-1)/2$ & প্রতি নোডে $n-1$ পোর্ট; মোট পোর্টপ্রান্ত $n(n-1)=2L$। প্রতি জোড়ায় একটি দ্বিমুখী লিংক ধরা হয়। Partial Mesh-এ একই সূত্র প্রযোজ্য নয়।\\\hline
\end{tabularx}
Bus-এর ভাগাভাগি করা ব্যাকবোন প্রতিটি কম্পিউটার-জোড়ার পৃথক লিংক নয়। টপোলজির নিজস্ব স্থির ডেটা হার বা কিলোমিটারে নির্দিষ্ট পরিসর থাকে না। লিংক, যন্ত্র ও নকশা এসব নির্ধারণ করে।

\subsection{Full Mesh-এর Link, Port ও Node-এর হিসাব}
\label{ch2:c09-08}
প্রতিটি নোড অন্য $n-1$টির সঙ্গে যুক্ত থাকায় $n(n-1)$ প্রান্ত পাওয়া যায়। এখানে প্রতিটি লিংকের দুই প্রান্ত গণনা হয়েছে। লিংকের সংখ্যা পেতে তাই ২ দিয়ে ভাগ করতে হয়। দুই দিকে ডেটা যায় বলে একই লিংক দুবার গণনা করা যাবে না। আগের ছকের সূত্রে কয়েকটি উদাহরণ নিচে দেওয়া হলো:
\chTwoPair{\textbf{উদাহরণ: ছয় সার্ভার}\par $L=6\times5/2=15$ লিংক। প্রতি সার্ভারে ৫ পোর্ট এবং মোট ৩০টি পোর্টপ্রান্ত। চিত্রে প্রতিটি জোড়া একবার যুক্ত। রেখা পরস্পরকে অতিক্রম করলেই সেখানে অতিরিক্ত নোড হয় না।}{\textbf{১২ নোডের হিসাব:} $12\times11/2=66$ লিংক; প্রতি নোডে ১১ পোর্ট।\par
\textbf{উল্টো হিসাব: ৪৫ লিংক}\par Full Mesh হলে $n(n-1)/2=45$, অর্থাৎ $n(n-1)=90$। $10\times9=90$, তাই ১০ নোড। প্রতি নোডে ৯ পোর্ট ও মোট ৯০টি পোর্টপ্রান্ত।}
\textbf{নতুন নোড যোগ:} পুরোনো $n$টি নোডের প্রত্যেকটির সঙ্গে নতুন নোড যুক্ত করতে $n$টি নতুন লিংক লাগে। নতুন নোডে $n$টি পোর্ট এবং প্রত্যেক পুরোনো নোডে একটি বাড়তি পোর্ট প্রয়োজন। যেমন, চার নোডের ৬ লিংক থেকে পাঁচ নোডের ১০ লিংক করতে আরও ৪ লিংক লাগে।
\textbf{কয়েকটি নোড একসঙ্গে যোগ:} $n$টি থেকে $n+k$টি নোড হলে বাড়তি লিংকের সংখ্যা
$k n+k(k-1)/2$। প্রথম অংশ নতুন ও পুরোনো নোডের জোড়া বোঝায়। দ্বিতীয় অংশ নতুন নোডগুলোর নিজেদের জোড়া বোঝায়। পাঁচ থেকে সাত নোড হলে $2\times5+2\times1/2=11$ নতুন লিংক লাগে; মোট $10+11=21$।

\subsection{Link বদলালে Topology-ও বদলায়}
\label{ch2:c09-09}
নোডের সংখ্যা একই রেখে লিংক যোগ বা বাদ দিলে সংযোগবিন্যাস বদলাতে পারে। চার নোড P, Q, R, S-এর Full Mesh-এ ৬টি লিংক আছে: PQ, QR, RS, SP, PR ও QS। PR এবং QS বাদ দিলে PQ–QR–RS–SP বন্ধ চক্র থাকে। প্রতিটি নোডের দুই প্রতিবেশী থাকায় অবশিষ্ট সংযোগ Ring বিন্যাস হয়।
\chTwoPair{\centering
\begin{tikzpicture}[x=1cm,y=1cm]
\foreach \x/\y/\l in {0/1.9/P,3/1.9/Q,3/0/R,0/0/S}{\coordinate (p\l) at (\x,\y);}
\foreach \a/\b in {P/Q,Q/R,R/S,S/P}{\draw[chTwoLink] (p\a)--(p\b);}
\draw[chTwoLink,draw=chTwoRose,dashed] (pP)--(pR);
\draw[chTwoLink,draw=chTwoRose,dashed] (pQ)--(pS);
\foreach \l in {P,Q,R,S}{\node[chTwoNode,circle] at (p\l){\l};}
\end{tikzpicture}\par\textbf{আগে:} ৪ Node, ৬ Link; Dashed দুটিও সংযোগ।}{\centering
\begin{tikzpicture}[x=1cm,y=1cm]
\foreach \x/\y/\l in {0/1.9/P,3/1.9/Q,3/0/R,0/0/S}{\coordinate (p\l) at (\x,\y);}
\foreach \a/\b in {P/Q,Q/R,R/S,S/P}{\draw[chTwoLink] (p\a)--(p\b);}
\foreach \l in {P,Q,R,S}{\node[chTwoNode,circle] at (p\l){\l};}
\end{tikzpicture}\par\textbf{পরে:} PR ও QS বাদ; ৪ Link-এর চক্র।}
যেকোনো দুটি লিংক বাদ দিলেই Ring হয় না। লিংক দুটি কোন নোড সংযুক্ত করেছে, তা বিবেচ্য। একই নোডের দুটি লিংক বাদ দিলে তার একটি মাত্র সংযোগ থাকতে পারে। তখন চার নোডের ওই চক্র থাকে না। রিংয়ে PR ও QS দুটিই যোগ করলে Full Mesh ফিরে আসে। একটি যোগ করলে সম্পূর্ণ Full Mesh হয় না। চিত্রের এই পরিবর্তন থেকে Token Passing বা ডেটা চলাচলের দিক নিজে থেকে নির্ধারণ করা যায় না।
